\documentclass[11pt]{article}
\usepackage[paperwidth=11in,paperheight=8.5in,margin=0.25in]{geometry}
\usepackage[T1]{fontenc}
\usepackage[utf8]{inputenc}
\usepackage[english]{babel}
\usepackage{lmodern}
\usepackage{microtype}
\usepackage{amsmath,amssymb}
\usepackage{xcolor}
\usepackage{tikz}
\usetikzlibrary{calc}
\hyphenpenalty=8000
\exhyphenpenalty=8000
\emergencystretch=2em
\definecolor{QCBlue}{HTML}{1F4E79}
\definecolor{QCGreen}{HTML}{2E6B5F}
\definecolor{QCOrange}{HTML}{B65B2A}
\definecolor{QCPurple}{HTML}{4B3F72}
\definecolor{QCLight}{HTML}{F4F6F8}
\definecolor{QCBorder}{HTML}{D0D6DC}
\definecolor{QCText}{HTML}{111111}
\newlength{\Margin} \setlength{\Margin}{0.25in}
\newlength{\Gap} \setlength{\Gap}{0.1in}
\newlength{\TitleH} \setlength{\TitleH}{0.8in}
\newlength{\HeaderH} \setlength{\HeaderH}{0.38in}
\newlength{\Pad} \setlength{\Pad}{0.1in}
\newlength{\UsableW}
\newlength{\CellW}
\newlength{\CellH}
\setlength{\UsableW}{\dimexpr\paperwidth-2\Margin\relax}
\setlength{\CellW}{5.2in}
\setlength{\CellH}{3.5in}
\newcommand{\QuadBox}[4]{%
  \node[anchor=north west, draw=QCBorder, line width=0.6pt, rounded corners=10pt, fill=white, minimum width=\CellW, minimum height=\CellH, inner sep=0pt] (B) at #1 {};
  \fill[#2, rounded corners=10pt] ([xshift=0pt,yshift=0pt]B.north west) rectangle ([xshift=0pt,yshift=-\HeaderH]B.north east);
  \node[anchor=west, text=white, font=\bfseries] at ([xshift=0.18in,yshift=-0.22in]B.north west) {#3};
  \node[anchor=north west, text=QCText] at ([xshift=\Pad,yshift=-\HeaderH-\Pad]B.north west) {\begin{minipage}[t]{\dimexpr\CellW-2\Pad\relax}\raggedright\small #4\end{minipage}};
}
\pagestyle{empty}
\begin{document}
\begin{tikzpicture}[remember picture,overlay]
\coordinate (NW) at ([xshift=\Margin,yshift=-\Margin]current page.north west);
\node[anchor=north west, draw=QCBorder, line width=0.6pt, rounded corners=12pt, fill=QCLight, minimum width=\UsableW, minimum height=\TitleH, inner sep=0pt] (T) at (NW) {};
\node[anchor=north west] at ([xshift=0.22in,yshift=-0.22in]T.north west) {\begin{minipage}[t]{\dimexpr\UsableW-0.44in\relax}{\Huge\bfseries \textcolor{QCBlue}{Dielectrics Formula Sheet}}\par\vspace{2pt}{\normalsize \textbf{Diego Juarez}\hfill \textbf{Computational Electrodynamics}\hfill \textbf{\today}}\end{minipage}};
\coordinate (GridNW) at ([yshift=-0.15in]T.south west);
\coordinate (TL) at (GridNW);
\coordinate (TR) at ([xshift=\CellW+\Gap]GridNW);
\coordinate (BL) at ([yshift=-\CellH-\Gap]GridNW);
\coordinate (BR) at ([xshift=\CellW+\Gap,yshift=-\CellH-\Gap]GridNW);
\QuadBox{(TL)}{QCBlue}{1) Polarization and D Field}{
\[
\mathbf P=\text{dipole moment per unit volume}
\]
Bound charges:
\[
\rho_b=-\nabla\cdot \mathbf P, \qquad \sigma_b=\mathbf P\cdot \hat{\mathbf n}
\]
Displacement field:
\[
\mathbf D=\varepsilon_0\mathbf E+\mathbf P
\]
Gauss law in matter:
\[
\nabla\cdot \mathbf D=\rho_f
\]
Linear dielectric:
\[
\mathbf P=\varepsilon_0\chi_e \mathbf E
\]
\[
\mathbf D=\varepsilon \mathbf E,\qquad \varepsilon=\varepsilon_0(1+\chi_e)
\]
}
\QuadBox{(TR)}{QCGreen}{2) Boundary Conditions}{
Tangential electric field:
\[
\mathbf E_{1t}=\mathbf E_{2t}
\]
Normal displacement:
\[
D_{2n}-D_{1n}=\sigma_f
\]
No free surface charge:
\[
D_{1n}=D_{2n}
\]
Potential form:
\[
\nabla\cdot(\varepsilon\nabla\Phi)=-\rho_f
\]
Constant $\varepsilon$:
\[
\nabla^2\Phi=-\frac{\rho_f}{\varepsilon}
\]
}
\QuadBox{(BL)}{QCOrange}{3) Canonical Results}{
Uniformly polarized sphere:
\[
\rho_b=0,\qquad \sigma_b=P_0\cos\theta
\]
Dielectric-filled capacitor:
\[
E=\frac{V}{d}, \qquad C=\varepsilon\frac{A}{d}
\]
Point charge in uniform dielectric:
\[
\Phi=\frac{q}{4\pi\varepsilon r}
\]
Dielectric sphere in uniform field:

Interior field is uniform and reduced relative to the applied field.
}
\QuadBox{(BR)}{QCPurple}{4) Quick Reminders}{
\textbf{Free charge} is externally supplied.

\textbf{Bound charge} arises from polarization.

\textbf{Useful checks}

1. Use $\mathbf D$ for free-charge bookkeeping.

2. Use $\mathbf P$ to find bound surface and volume charge.

3. At interfaces, do not force $\mathbf E_n$ to be continuous.

4. In matter, the relevant equation is often $\nabla\cdot(\varepsilon\nabla\Phi)=-\rho_f$.
}
\end{tikzpicture}
\end{document}
